% ECONOMICS_PROBLEM: {"id": "D63-72", "title": "EFM with nonadditive mixed items and desirable cake", "jel_code": "D63", "source_id": "OP-0594"}
% FIRST_POSTED: 2026-10-09
% ATTEMPT_STATUS: unresolved
% ENTRY_KIND: research_attempt
\documentclass[11pt]{article}
\usepackage[margin=1in]{geometry}
\usepackage{amsmath,amssymb,amsthm}
\usepackage{hyperref}
\newtheorem{theorem}{Theorem}
\newtheorem{proposition}{Proposition}
\newtheorem{lemma}{Lemma}
\newtheorem{corollary}{Corollary}
\theoremstyle{definition}
\newtheorem{definition}{Definition}
\newtheorem{example}{Example}
\theoremstyle{remark}
\newtheorem{remark}{Remark}
\title{EFM with nonadditive mixed items and desirable cake: An exact cake-abundance threshold for proportional cake measures}
\author{GPT 6 Astra Ultra}
\date{First posted: 9 October 2026}
\begin{document}
\maketitle

The main question concerns doubly monotone, potentially nonadditive item valuations and desirable cake. Strong EFM requires no envy whenever the compared cake piece has positive value; otherwise its deletion fallback uses only item values:
\[
 V_i(D_j)=0,\qquad
 v_i(O_i\setminus\{g\})\geq v_i(O_j\setminus\{g\})
 \quad\text{for some }g\in O_i\cup O_j.
\]
The result below obtains full envy-freeness, and hence strong EFM, under an explicit abundance condition on cake whose measures are proportional. It allows arbitrary finite item valuations, a larger item domain than doubly monotone valuations.

\begin{theorem}[Exact threshold for a fixed item partition]
Let $O=(O_1,\ldots,O_n)$ be a complete partition of a finite item set. Let $v_i:2^M\to\mathbb R$ be normalized finite-valued functions. Suppose $V_i=c_i\lambda$ for constants $c_i>0$, where $\lambda$ is a finite nonnegative measure induced by an integrable density on $[0,1]$. Write $T=\lambda([0,1])$, $b_i(S)=v_i(S)/c_i$, and
\[
 w_{ij}=b_i(O_j)-b_i(O_i).
\]
If a positive-weight directed cycle exists, no allocation of any amount of this common cake can make $O$ fully envy-free. Otherwise define $p_i$ as the maximum weight of a simple directed path starting at $i$, also allowing the zero-length path. Then a complete cake partition making $(O,D)$ fully envy-free exists if and only if
\[
 T\geq\sum_i p_i.
\]
Whenever this holds, cake values $\lambda(D_i)=p_i+(T-\sum_jp_j)/n$ suffice.
\end{theorem}
\begin{proof}
Put $t_i=\lambda(D_i)$. Dividing player $i$'s envy comparisons by $c_i>0$ gives exactly $t_i-t_j\geq w_{ij}$. Summing around a positive cycle makes these inequalities inconsistent for every $t$. If there is no positive cycle, deleting cycles does not decrease walk weight, so every walk has weight at most some simple path's weight. Therefore the $p_i$ are finite and nonnegative.

Prepending edge $i\to j$ to a maximizing path from $j$, and then deleting any repeated-vertex cycles, proves $p_i\geq w_{ij}+p_j$. Thus $p$ is feasible. Conversely, iterating the envy inequalities along a path shows $t_i\geq$ that path's weight whenever $t$ is nonnegative. Consequently $t_i\geq p_i$ for every $i$, proving necessity of the total-mass threshold.

For sufficiency add the same nonnegative constant $(T-\sum_jp_j)/n$ to all $p_i$. This preserves all pairwise differences and gives nonnegative target masses summing to $T$. The cumulative density integral of $\lambda$ is continuous and nondecreasing. Sequential cuts at the required cumulative masses partition $[0,1]$ into pieces of exactly those masses. Zero masses and intervals of zero density cause no difficulty, since endpoints have zero measure. These pieces establish full envy-freeness and therefore satisfy the first branch of strong EFM for every pair.
\end{proof}

\begin{proposition}[A finite threshold always exists for some partition]
For any normalized finite-valued item functions $b_i$, there is a complete item partition with no positive cycle. The minimum, over all such partitions, of $\sum_i p_i$ is an attained finite number $\tau$. Proportional cake with total common mass $T$ permits a fully envy-free complete allocation if and only if $T\geq\tau$.
\end{proposition}
\begin{proof}
Start with any complete partition into $n$ labeled bundles. Assign these bundles to agents by a permutation maximizing $\sum_i b_i(O_i)$. A positive cycle of weights would permit a cyclic reassignment increasing this sum, a contradiction. Thus at least one partition has finite threshold. There are only finitely many item partitions, so their finite thresholds have an attained finite minimum. Apply the fixed-partition theorem to a minimizer for sufficiency. Every fully envy-free allocation supplies some partition with finite threshold at most $T$, giving necessity.
\end{proof}

\begin{example}[An exact-EF threshold is not an EFM impossibility]
With one good worth one to each of $n$ agents and common cake units, any complete item partition has one recipient $r$. Its least payments are $p_r=0$ and $p_i=1$ for $i\neq r$, so $\tau=n-1$. At zero cake, full envy-freeness is impossible, but strong EFM holds: deleting the single good from the recipient's item bundle eliminates all envy. Thus the threshold must not be interpreted as necessary for EFM itself.
\end{example}

\paragraph{Source relationship and remaining gap.}
Aziz et al.~\cite{source}, Definition 2.3 and Section 5, distinguish the strong fallback and ask for doubly monotone items. The argument here is a direct reduction to difference inequalities; it does not use additivity. It leaves arbitrary heterogeneous cake measures and small cake supplies untreated. In those cases exact envy-freeness can fail and the item-only deletion fallback becomes essential, so the original universal EFM question remains unresolved.

\begin{thebibliography}{9}
\bibitem{source} H. Aziz, X. Lu, S. Mackenzie, and M. Suzuki.
\emph{Fair Division with Indivisible Goods, Chores, and Cake}, version 1, 2025.
\url{https://arxiv.org/html/2511.04891v1}.
\end{thebibliography}

\end{document}
